\documentclass[12pt,a4paper]{article}
\usepackage[utf8]{inputenc}
\usepackage[vietnamese]{babel}
\usepackage{amsmath,amssymb,amsthm}
\usepackage{geometry}
\usepackage{enumitem}
\usepackage[most]{tcolorbox}
\usepackage{hyperref}
\usepackage{fancyhdr}
\usepackage{tikz}
\usepackage{eso-pic}
\usepackage{xcolor}
\geometry{a4paper, margin=2cm, bottom=2.5cm}
\hypersetup{colorlinks=true, linkcolor=blue!70!black}
\setlist[itemize]{topsep=4pt, itemsep=2pt}
\setlist[enumerate]{topsep=4pt, itemsep=2pt}
\AddToShipoutPictureFG{%
  \AtPageCenter{%
    \begin{tikzpicture}[remember picture, overlay]
      \node[rotate=45, scale=1, opacity=0.06]{\fontsize{60}{72}\selectfont\bfseries\sffamily Đặng Tấn Quí};
    \end{tikzpicture}%
  }%
}
\pagestyle{fancy}
\fancyhf{}
\fancyhead[L]{\small\textit{TOÁN HỌC TÍNH TOÁN}}
\fancyhead[R]{\small\textit{Đặng Tấn Quí}}
\fancyfoot[C]{\thepage}
\fancyfoot[L]{\footnotesize\textcopyright\ Đặng Tấn Quí}
\fancyfoot[R]{\footnotesize SĐT: 0377\,122\,210}
\renewcommand{\headrulewidth}{0.4pt}
\renewcommand{\footrulewidth}{0.4pt}
\fancypagestyle{plain}{\fancyhf{}\fancyfoot[C]{\thepage}\fancyfoot[L]{\footnotesize\textcopyright\ Đặng Tấn Quí}\fancyfoot[R]{\footnotesize SĐT: 0377\,122\,210}\renewcommand{\headrulewidth}{0pt}\renewcommand{\footrulewidth}{0.4pt}}
\newtcolorbox{dinhnghia}[1][]{enhanced, breakable, colback=blue!3!white, colframe=blue!60!black, fonttitle=\bfseries, title={Định nghĩa}, #1}
\newtcolorbox{dinhly}[1][]{enhanced, breakable, colback=green!3!white, colframe=green!50!black, fonttitle=\bfseries, title={Định lý}, #1}
\newtcolorbox{tinhchat}[1][]{enhanced, breakable, colback=yellow!5!white, colframe=orange!50!black, fonttitle=\bfseries, title={Tính chất}, #1}
\newtcolorbox{phuongphap}[1][]{enhanced, breakable, colback=gray!5!white, colframe=gray!50!black, fonttitle=\bfseries, title={Phương pháp}, #1}
\newtcolorbox{luuy}[1][]{enhanced, breakable, colback=red!3!white, colframe=red!50!black, fonttitle=\bfseries, title={Lưu ý}, #1}
\newtcolorbox{congthuc}[1][]{enhanced, breakable, colback=cyan!3!white, colframe=cyan!50!black, fonttitle=\bfseries, title={Công thức}, #1}
\newtcolorbox{vidu}[1][]{enhanced, breakable, colback=violet!3!white, colframe=violet!50!black, fonttitle=\bfseries, title={Ví dụ}, #1}

\begin{document}
\thispagestyle{empty}
\begin{tikzpicture}[remember picture, overlay]
  \fill[blue!80!black] (current page.south west) rectangle (current page.north east);
  \node[anchor=west, white, font=\fontsize{26}{32}\bfseries\selectfont]
    at ([xshift=3cm, yshift=-7.5cm]current page.north west) {TOÁN HỌC TÍNH TOÁN};
  \node[anchor=west, white, font=\fontsize{18}{24}\selectfont]
    at ([xshift=3cm, yshift=-9cm]current page.north west) {Độ phức tạp, P vs NP, thuật toán số, Turing};
  \node[anchor=west, white, font=\Large\bfseries] at ([xshift=3cm, yshift=-13cm]current page.north west) {Biên soạn:};
  \node[anchor=west, yellow!80!white, font=\fontsize{22}{28}\bfseries\selectfont] at ([xshift=3cm, yshift=-14.5cm]current page.north west) {ĐẶNG TẤN QUÍ};
  \node[anchor=south east, white, font=\Large\bfseries] at ([xshift=-2cm, yshift=3cm]current page.south east) {\the\year};
\end{tikzpicture}
\newpage
\tableofcontents
\newpage
%% ================================================================
%%  CHƯƠNG 1: ĐỘ PHỨC TẠP TÍNH TOÁN
%% ================================================================
\section{Độ phức tạp tính toán}

\begin{dinhnghia}[title={Ký hiệu tiệm cận}]
  \begin{itemize}
    \item $f = O(g)$: $\exists c, n_0$: $|f(n)| \le c\,g(n)$ $\forall n \ge n_0$ (cận trên).
    \item $f = \Omega(g)$: $\exists c, n_0$: $f(n) \ge c\,g(n)$ $\forall n \ge n_0$ (cận dưới).
    \item $f = \Theta(g)$: $f = O(g)$ và $f = \Omega(g)$.
    \item $f = o(g)$: $f/g \to 0$.
  \end{itemize}
\end{dinhnghia}

\begin{dinhnghia}[title={Mô hình tính toán --- Máy Turing}]
  $(Q, \Sigma, \Gamma, \delta, q_0, q_{\text{acc}}, q_{\text{rej}})$: tập trạng thái, bảng chữ cái input, bảng chữ băng, hàm chuyển, trạng thái đầu/chấp nhận/từ chối.

  $L(M)$: ngôn ngữ mà $M$ chấp nhận.
\end{dinhnghia}

\begin{dinhnghia}[title={Lớp phức tạp}]
  \begin{itemize}
    \item \textbf{P}: các bài toán quyết định giải được trong thời gian đa thức (deterministc).
    \item \textbf{NP}: xác minh lời giải trong thời gian đa thức (hoặc: máy Turing không tất định đa thức).
    \item \textbf{co-NP}: phần bù của NP.
    \item \textbf{PSPACE}: giải được với bộ nhớ đa thức.
  \end{itemize}
\end{dinhnghia}

\begin{dinhnghia}[title={Quy dẫn và NP-đầy đủ}]
  $A \le_p B$ ($A$ quy dẫn đa thức về $B$): $\exists f$ đa thức sao cho $x \in A \Leftrightarrow f(x) \in B$.

  $B$ \textbf{NP-khó}: mọi $A \in \text{NP}$ đều $A \le_p B$. $B$ \textbf{NP-đầy đủ}: NP-khó $\cap$ NP.
\end{dinhnghia}

\begin{dinhly}[title={Cook--Levin}]
  SAT (bài toán thỏa mãn) là NP-đầy đủ.
\end{dinhly}

\begin{vidu}[title={Bài toán NP-đầy đủ quan trọng}]
  SAT, 3-SAT, Clique, Vertex Cover, Hamilton Cycle, TSP (quyết định), Graph Coloring, Subset Sum, Knapsack, \ldots
\end{vidu}

\begin{luuy}
  Câu hỏi mở nổi tiếng: $\textbf{P} \overset{?}{=} \textbf{NP}$ (bài toán triệu đô Clay).
\end{luuy}

%% ================================================================
%%  CHƯƠNG 2: THUẬT TOÁN SỐ HỌC
%% ================================================================
\section{Thuật toán số học}

\begin{phuongphap}[title={Lũy thừa modular nhanh}]
  Tính $a^n \bmod m$ bằng bình phương liên tiếp: $O(\log n)$ phép nhân modular.
\end{phuongphap}

\begin{phuongphap}[title={Kiểm tra nguyên tố}]
  \begin{itemize}
    \item \textbf{Chia thử}: $O(\sqrt{n})$.
    \item \textbf{Miller--Rabin}: xác suất, $O(k\log^2 n)$ ($k$ vòng).
    \item \textbf{AKS}: tất định đa thức (lý thuyết), $O(\log^{6+\varepsilon} n)$.
  \end{itemize}
\end{phuongphap}

\begin{phuongphap}[title={Phân tích thừa số}]
  \begin{itemize}
    \item \textbf{Chia thử}: $O(\sqrt{n})$.
    \item \textbf{Pollard's $\rho$}: $O(n^{1/4})$ heuristic.
    \item \textbf{Sàng trường số (NFS)}: $O\!\left(\exp\!\left(c(\ln n)^{1/3}(\ln\ln n)^{2/3}\right)\right)$ --- tốt nhất hiện nay.
  \end{itemize}
\end{phuongphap}

\begin{phuongphap}[title={Logarit rời rạc}]
  Tìm $x$: $g^x \equiv h \pmod{p}$.
  \begin{itemize}
    \item \textbf{Baby-step Giant-step} (Shanks): $O(\sqrt{p})$ thời gian và bộ nhớ.
    \item \textbf{Pollard's $\rho$}: $O(\sqrt{p})$ thời gian, $O(1)$ bộ nhớ.
    \item \textbf{Index Calculus}: dưới mũ (cho $\mathbb{F}_p^*$), không áp dụng cho ECC.
  \end{itemize}
\end{phuongphap}

%% ================================================================
%%  CHƯƠNG 3: TÍNH TOÁN VÀ QUYẾT ĐỊNH ĐƯỢC
%% ================================================================
\section{Tính toán và quyết định được}

\begin{dinhnghia}[title={Quyết định được}]
  Ngôn ngữ $L$ \textbf{quyết định được} (decidable/recursive) nếu tồn tại máy Turing luôn dừng và chấp nhận đúng $L$.
\end{dinhnghia}

\begin{dinhnghia}[title={Liệt kê được}]
  $L$ \textbf{liệt kê được} (recursively enumerable, r.e.) nếu tồn tại MT chấp nhận mọi $x \in L$ (có thể không dừng với $x \notin L$).
\end{dinhnghia}

\begin{dinhly}[title={Bài toán dừng (Halting Problem)}]
  Không tồn tại MT quyết định: ``MT $M$ có dừng trên input $w$?''. Bài toán dừng là \textbf{không quyết định được} nhưng liệt kê được.
\end{dinhly}

\begin{dinhly}[title={Định lý Rice}]
  Mọi tính chất không tầm thường của ngôn ngữ r.e. đều không quyết định được.
\end{dinhly}

\end{document}